% ============================================================
% Chapitre 7 — Conclusion générale
% ============================================================

\chapter{Conclusion générale}
\label{ch:conclusion}

\section*{Bilan du travail}
\phantomsection
\addcontentsline{toc}{section}{Bilan du travail}

Ce mémoire avait pour objectif de concevoir et de
valider un système de gestion intelligente des feux
de circulation, déployable sur Raspberry Pi 5, capable
de détecter les usagers de la route, d'anonymiser les
piétons et de produire une décision de signalisation
adaptative. La question centrale était donc de savoir
s'il est possible de combiner \textbf{Edge AI},
\textbf{Privacy-by-Design} et \textbf{décision MDP}
dans un pipeline unique, léger et exploitable dans un
contexte de smart city marocaine.

Les résultats obtenus montrent que cette approche est
réalisable. Le pipeline final intègre six modules :
acquisition vidéo, détection YOLO26n, anonymisation,
tracking par zones polygonales, décision MDP et
dashboard Flask. L'ensemble fonctionne sans dépendance
cloud et respecte une architecture compatible avec les
contraintes matérielles d'un dispositif embarqué.

\section*{Résultats principaux}
\phantomsection
\addcontentsline{toc}{section}{Résultats principaux}

Le modèle final \textbf{YOLO26n Step 3 960} atteint
\textbf{0.792 mAP50} et \textbf{0.579 mAP50-95} sur
le split test. Les classes véhicules les plus importantes
pour la décision de trafic obtiennent de très bons scores,
notamment \textit{car} (0.931 mAP50),
\textit{motorcycle} (0.915 mAP50) et
\textit{emergency\_vehicle} (0.926 mAP50). La classe
\textit{person} progresse par rapport au modèle 800,
mais reste la limite principale du système avec un rappel
test de 0.353.

Le benchmark embarqué confirme que la famille YOLO26n
nano est compatible avec le Raspberry Pi 5. La configuration
FP32 benchmarkée atteint \textbf{7.6 FPS réels}, au-dessus
du seuil minimal de 5 FPS retenu pour la gestion adaptative
des feux. Le modèle 960 constitue le meilleur choix de
précision; l'export 800 améliore le débit sur Pi jusqu'à
\textbf{3.94 FPS} lorsque chaque frame est inférée.
Avec \texttt{--vid-stride 2}, le même export couvre
\textbf{7.84 frames source/s}, au prix d'une inférence
sur une frame source sur deux. La quantification INT8
dynamique a été testée, mais elle n'est pas retenue car
elle réduit la taille du modèle tout en descendant à
\textbf{2.06 FPS} sur CPU ARM.

Le module d'anonymisation applique un floutage gaussien
sur toutes les boîtes détectées de classe \textit{person}
avant le tracking. Cette position dans le pipeline réduit
la circulation de pixels identifiants et concrétise le
principe Privacy-by-Design. Enfin, la politique MDP produit
la décision attendue sur les cinq scénarios de validation :
trafic dense Nord-Sud, trafic Est-Ouest, urgence,
priorité piétonne et absence de trafic.

\begin{table}[H]
\centering
\caption{Synthèse finale des hypothèses}
\small
\begin{tabularx}{\textwidth}{clX}
\toprule
\textbf{H} & \textbf{Axe} & \textbf{Conclusion} \\
\midrule
H1 & Détection &
Validée avec une limite sur les petits piétons :
0.792 mAP50 test, rappel \textit{person} de 0.353. \\
H2 & Dataset hybride &
Validée empiriquement : le passage 800 vers 960 améliore
la mAP50 globale et la classe \textit{person}. \\
H3 & Edge AI &
Validée pour la famille YOLO26n nano à 7.6 FPS sur Pi 5;
Step 3 atteint 7.84 frames source/s avec l'export 800
et \texttt{--vid-stride 2}, mais seulement une frame
sur deux est inférée. \\
H4 & Privacy-by-Design &
Validée sur les détections : 100\% des boîtes
\textit{person} détectées sont floutées avant tracking. \\
H5 & Décision MDP &
Validée sur cinq scénarios représentatifs. \\
\bottomrule
\end{tabularx}
\label{tab:conclusion_hypotheses}
\end{table}

\section*{Limites}
\phantomsection
\addcontentsline{toc}{section}{Limites}

La première limite concerne la détection des petits piétons.
Même avec CityPersons, les tiny crops, une résolution de
960 pixels et les ROI, certains piétons éloignés ou occultés
peuvent ne pas être détectés. Cette limite entraîne un risque
résiduel pour l'anonymisation, car seules les personnes
détectées peuvent être floutées.

La seconde limite concerne le compromis précision/latence sur
Raspberry Pi 5. Les résultats Pi valident la famille YOLO26n
nano en temps réel, mais le pipeline complet Step 3 960 avec
ROI atteint 2.16 FPS et monte seulement à 2.51 FPS sans ROI
additionnelles. L'export 800 monte à 3.94 FPS sans ROI et
3.45 FPS avec ROI personnes lorsque toutes les frames
sont inférées. Le mode \texttt{--vid-stride 2} atteint
7.84 frames source/s, ce qui valide un compromis de
couverture vidéo, mais pas une inférence exhaustive de
chaque frame. La quantification INT8 dynamique a été
écartée car elle ralentit le modèle Step 3 à
2.06 FPS sur Pi 5.
Enfin, le dataset local
marocain reste encore limité : un déploiement réel à Agadir
nécessiterait davantage de vidéos couvrant la nuit, la pluie,
les intersections complexes et les véhicules atypiques.

\section*{Perspectives}
\phantomsection
\addcontentsline{toc}{section}{Perspectives}

Les perspectives prioritaires sont l'élargissement du dataset
marocain, l'amélioration de la classe \textit{person}, le test
de stratégies d'accélération matérielle ou de planification
d'inférence adaptative, et l'étude d'une coordination
multi-intersections. À moyen terme, la politique
MDP déterministe pourrait être comparée à un agent de
\textit{Deep Reinforcement Learning} distillé pour edge device.

Ce travail montre qu'une solution locale, explicable et
respectueuse des données personnelles peut constituer une
base réaliste pour des expérimentations smart city à Agadir.
Il ne remplace pas une infrastructure urbaine complète, mais
il démontre qu'un prototype intelligent, embarqué et conforme
aux contraintes juridiques peut être construit avec un coût
matériel réduit et une architecture maîtrisable.
